\section{Elemen Semantik dan Struktur Halaman}
Elemen semantik membantu memberi makna pada struktur halaman, misalnya \code{header}, \code{nav}, \code{main}, dan \code{footer}. Dengan elemen semantik, mesin pencari dan pembaca layar lebih mudah memahami konten \cite{webdevHtml}. Ini juga memudahkan pengembang lain memahami struktur halaman.

Struktur semantik dapat dianggap sebagai kerangka informasi halaman. Misalnya, \code{header} berisi judul dan navigasi, sedangkan \code{main} berisi konten utama. Kebiasaan menggunakan elemen semantik sejak awal akan meningkatkan kualitas proyek Anda.

\begin{table}[ht]
\centering
\begin{tabular}{|P{4cm}|P{7cm}|}
\hline
\textbf{Elemen} & \textbf{Fungsi} \\
\hline
\code{header} & Bagian atas halaman, judul dan navigasi \\
\hline
\code{main} & Konten utama halaman \\
\hline
\code{section} & Pengelompokan konten terkait \\
\hline
\code{footer} & Informasi penutup, hak cipta \\
\hline
\end{tabular}
\caption{Contoh elemen semantik HTML5}
\end{table}

